\section{Utilisation de BerTopic}

\subsection{Lemmatisation des données}

Comme pour tout modèle de traitement du langage naturel, il est important de préparer les données avant de les utiliser avec BerTopic. La lemmatisation est une étape clé dans ce processus, j'ai suivi les memes etapes que pour le NMF ainsi que le KMeans. 

J'ai utilisé la bibliothèque spaCy pour effectuer la lemmatisation des textes, j'ai réutilisé le même modele de langue que pour le \textit{Topic modeling} classique, le \textit{"fr\_core\_news\_lg"}. Cette étape permet de réduire les mots à leur forme de base, ce qui facilite l'analyse et l'extraction des thèmes. 

\subsection{Préparation des embeddings}

Le modèle \texttt{dangvantuan/sentence-camembert-base} est utilisé afin de transformer chaque segment textuel en un vecteur numérique dense, appelé \textit{embedding}, qui représente son contenu sémantique. Fondé sur CamemBERT et adapté à la langue française, ce modèle permet de rapprocher les documents qui traitent de thèmes similaires, même lorsqu’ils n’emploient pas exactement le même vocabulaire. Ces représentations vectorielles sont ensuite fournies à BERTopic, qui s’appuie sur leur proximité pour réduire leur dimension, regrouper les documents sémantiquement proches et identifier les principaux topics du corpus. L’utilisation de ces embeddings est particulièrement pertinente dans le cadre d’un \textit{Dynamic Topic Modeling} (DTM), puisqu’elle permet d’étudier l’évolution des thèmes au cours du récit tout en tenant compte du contexte lexical et sémantique de chaque paquet de phrases.

\subsection{Préparation des paramètres du modèle}

Afin d’identifier des thèmes cohérents à partir des représentations sémantiques, BERTopic combine une réduction de dimension avec un algorithme de classification non supervisée. Le modèle UMAP réduit les embeddings à trois dimensions afin de faciliter la détection des groupes tout en conservant au mieux leur structure sémantique. Le paramètre \texttt{n\_neighbors=20} permet de prendre en compte un voisinage suffisamment large pour préserver à la fois les relations locales et l’organisation générale du corpus, tandis que \texttt{min\_dist=0.0} favorise la formation de groupes compacts. L’utilisation de la distance cosinus est adaptée aux embeddings textuels, car elle mesure la proximité entre leurs orientations plutôt que leur amplitude. Enfin, la fixation de \texttt{random\_state=42} garantit la reproductibilité de la réduction dimensionnelle.

La classification des segments est ensuite réalisée avec HDBSCAN. Le paramètre \texttt{min\_cluster\_size=27} impose qu’un topic soit représenté par au moins 27 segments, ce qui limite l’apparition de thèmes trop spécifiques ou insuffisamment documentés. Le paramètre \texttt{min\_samples=4} conserve une certaine souplesse dans la formation des clusters, tout en permettant d’identifier comme atypiques les observations éloignées des groupes principaux. La méthode \texttt{eom} privilégie les clusters les plus stables dans la structure hiérarchique produite par HDBSCAN. L’activation de \texttt{prediction\_data} conserve par ailleurs les informations nécessaires à l’affectation ultérieure de nouveaux documents ou au traitement des observations considérées comme aberrantes.

La représentation lexicale de chaque topic repose sur un \texttt{CountVectorizer}. Le seuil \texttt{min\_df=2} élimine les termes n’apparaissant que dans un seul segment, souvent peu représentatifs ou associés à des erreurs de transcription, tandis que \texttt{max\_df=0.6} écarte les mots présents dans plus de 60\,\% des segments, dont le pouvoir discriminant est généralement limité. Le modèle \texttt{ClassTfidfTransformer} pondère ensuite les termes selon leur importance dans chaque topic. L’option \texttt{reduce\_frequent\_words=True} réduit l’influence des mots fréquents, tandis que la pondération BM25 améliore la distinction entre les termes réellement caractéristiques d’un thème et ceux qui sont simplement courants dans le corpus.
Le modèle BERTopic est enfin ajusté sur les segments lemmatisés en utilisant les embeddings précédemment calculés. La désactivation du calcul des probabilités réduit le coût computationnel, celles-ci n’étant pas indispensables à l’identification et à l’interprétation des topics. Après l’entraînement, les segments initialement classés comme observations aberrantes sont réaffectés au topic sémantiquement le plus proche à partir de leurs embeddings. Cette étape de réduction des \textit{outliers}, suivie d’une mise à jour des représentations thématiques, permet d’obtenir une partition plus complète et des topics lexicalement plus cohérents, ce qui facilite ensuite l’étude de leur évolution dans le cadre du \textit{Dynamic Topic Modeling}.